\documentclass[11pt,a4paper]{article}
\usepackage[utf8]{inputenc}
\usepackage[T1]{fontenc}
\usepackage[spanish,es-minimal,es-noshorthands]{babel}
\AtBeginDocument{\renewcommand{\sin}{\operatorname{sen}}\renewcommand{\arcsin}{\operatorname{arcsen}}}
\usepackage{lmodern}
\usepackage[sfdefault,lining]{FiraSans}
\usepackage[a4paper,top=2.3cm,bottom=2.2cm,left=1.8cm,right=1.8cm,headheight=1.3cm,headsep=0.4cm,footskip=1cm]{geometry}
\usepackage{amsmath,amssymb}
\usepackage{icomma}
\usepackage{graphicx}
\usepackage[table]{xcolor}
\usepackage{tikz}
\usepackage{fontawesome5}
\usepackage{ccicons}
\usepackage{fancyhdr}
\usepackage{lastpage}
\usepackage{tasks}
\usepackage{needspace}
\usepackage[most]{tcolorbox}
\usepackage{array}
\usepackage[hidelinks]{hyperref}
\AtBeginDocument{\hypersetup{pdftitle={Soluciones de la ficha de la unidad 11: Probabilidad. 4º ESO},pdfsubject={Matemáticas A 4º ESO},pdfkeywords={probabilidad, experimento aleatorio, espacio muestral, sucesos, unión, intersección, suceso contrario, frecuencia relativa, ley de Laplace, experiencias compuestas, experiencias independientes, experiencias dependientes, probabilidad condicionada, diagrama de árbol, tabla de contingencia}}}

\definecolor{navy}{HTML}{1B2748}
\definecolor{accent}{HTML}{E8603A}
\definecolor{teal}{HTML}{1F9E90}
\definecolor{soft}{HTML}{EEF1F8}
\definecolor{muted}{HTML}{5D6685}

\setlength{\parindent}{0pt}
\setlength{\parskip}{3pt}
\settasks{label=\alph*), label-format=\bfseries\color{navy}, label-width=1.3em, item-indent=1.8em, column-sep=0.6em, after-item-skip=5pt, before-skip=3pt, after-skip=4pt}

% ---------- capçalera i peu ----------
\newcommand{\llicencia}{\href{https://creativecommons.org/licenses/by-nc-sa/4.0/deed.es}{\ccbyncsa}\ \ Licencia Creative Commons Atribución-NoComercial-CompartirIgual 4.0 Internacional (CC BY-NC-SA 4.0)}
\pagestyle{fancy}
\fancyhf{}
\fancyhead[L]{\raisebox{-0.35\height}{\includegraphics[height=0.95cm]{logo-gva.png}}}
\fancyhead[C]{\small\color{muted}Soluciones de la ficha · Unidad 11: Probabilidad · 4\textsuperscript{o} ESO}
\fancyhead[R]{\raisebox{-0.35\height}{\includegraphics[height=0.95cm]{logo-vermellar.png}}}
\fancyfoot[L]{\scriptsize\color{muted}\llicencia}
\fancyfoot[R]{\small\color{muted}\thepage}
\renewcommand{\headrulewidth}{0.4pt}
\renewcommand{\headrule}{\color{muted!40}\hrule height \headrulewidth}
\fancypagestyle{primera}{\fancyhf{}\renewcommand{\headrulewidth}{0pt}
  \fancyfoot[L]{\scriptsize\color{muted}\llicencia}
  \fancyfoot[R]{\small\color{muted}\thepage}}

% ---------- nivells i exercicis ----------
\newcommand{\estrelles}[1]{{\color{accent}\ifcase#1\or\faStar\faStar[regular]\faStar[regular]\or\faStar\faStar\faStar[regular]\or\faStar\faStar\faStar\fi}}
\newcounter{exer}
\newcommand{\ex}[1]{\par\needspace{4\baselineskip}\medskip\stepcounter{exer}\noindent{\color{navy}\bfseries\theexer.}\ \ {\small\estrelles{#1}}\ \ }

\newcommand{\seccio}[2]{\par\needspace{10\baselineskip}\bigskip
  \begin{tcolorbox}[colback=navy,colframe=navy,coltext=white,arc=2mm,boxrule=0pt,left=3mm,right=3mm,top=1.2mm,bottom=1.2mm]
  {\large\bfseries #1\quad #2}
  \end{tcolorbox}\vspace{-2pt}}

\newcolumntype{C}[1]{>{\centering\arraybackslash}p{#1}}
\usepackage{pifont}
\usetikzlibrary{trees}
% quadre de recordatori (només a la fitxa d'exercicis)
\newcommand{\recorda}[1]{\par\smallskip\begin{tcolorbox}[colback=accent!7,colframe=accent!55,arc=2mm,boxrule=0.5pt,left=2.5mm,right=2.5mm,top=1mm,bottom=1mm]\small{\color{accent}\faLightbulb[regular]}\ \textbf{Recuerda:}\ #1\end{tcolorbox}}
\newcommand{\res}[1]{\boxed{#1}}\newcommand{\rest}[1]{\fbox{#1}}

\def\fmneed{14}
% =====================================================================
% Capçalera FentMates de les fitxes: un color per a cada curs.
% Els colors són els de la pàgina inicial (LEVELS a assets/js/materials.js);
% el logo ve de _eines/logo/pdf/ (python3 _eines/logo/genera.py).
%
% Ús, al final del preàmbul de la unitat:
%   % =====================================================================
% Capçalera FentMates de les fitxes: un color per a cada curs.
% Els colors són els de la pàgina inicial (LEVELS a assets/js/materials.js);
% el logo ve de _eines/logo/pdf/ (python3 _eines/logo/genera.py).
%
% Ús, al final del preàmbul de la unitat:
%   % =====================================================================
% Capçalera FentMates de les fitxes: un color per a cada curs.
% Els colors són els de la pàgina inicial (LEVELS a assets/js/materials.js);
% el logo ve de _eines/logo/pdf/ (python3 _eines/logo/genera.py).
%
% Ús, al final del preàmbul de la unitat:
%   \input{../comu/fentmates.tex}
%   \fmcurs{1eso}
% i, al començament del document:
%   \capcalera{Matemàtiques}{1r ESO}{Fitxa d'exercicis · Unitat 2}{Divisibilitat}
% Opcional abans de \input: \def\fmneed{14} (espai mínim abans de cada secció).
% =====================================================================
\makeatletter
\def\fm@logodir{../../logo/pdf}
\providecommand{\fmneed}{10}
\def\fm@cursnom{}

% clau del curs, color i color del text damunt del color
\newcommand{\fm@curs}[3]{\@namedef{fm@color@#1}{#2}\@namedef{fm@text@#1}{#3}}
\fm@curs{1eso}{E6B84A}{navy}
\fm@curs{2eso}{3B82F6}{white}
\fm@curs{3eso}{FF7A50}{white}
\fm@curs{4eso-a}{A78BFA}{white}
\fm@curs{4eso-b}{6D28D9}{white}
\fm@curs{1batx-mi}{9B1C1C}{white}
\fm@curs{1batx-ccss}{EF6B6B}{white}
\fm@curs{2batx-mii}{1E6B47}{white}
\fm@curs{2batx-ccss}{2BB3A3}{white}

\newcommand{\fmcurs}[1]{%
  \@ifundefined{fm@color@#1}{\PackageError{fentmates}{Curs desconegut: #1}{}}{}%
  \edef\fm@tmp{\noexpand\definecolor{curs}{HTML}{\@nameuse{fm@color@#1}}%
               \noexpand\colorlet{curstext}{\@nameuse{fm@text@#1}}}\fm@tmp
  \colorlet{cursfosc}{curs!72!black}%
  \def\fm@logo{\fm@logodir/vertical-#1.pdf}%
  \def\fm@marca{\fm@logodir/marca-#1.pdf}}

% ---------- capçalera de la primera pàgina ----------
% #1 matèria (petit), #2 curs, #3 tipus de document i unitat, #4 títol
\newcommand{\capcalera}[4]{%
  \gdef\fm@cursnom{#2}%
  \par\noindent
  \begin{tcolorbox}[colback=curs,colframe=curs,boxrule=0pt,arc=2.2mm,width=3.5cm,height=2.15cm,
      valign=center,left=3mm,right=2mm,top=1mm,bottom=1mm,nobeforeafter,box align=center]
    \color{curstext}{\footnotesize #1}\par\vspace{1pt}{\fontsize{19}{21}\selectfont\bfseries #2\par}%
  \end{tcolorbox}\hspace{5mm}%
  \begin{minipage}[c]{\dimexpr\linewidth-3.5cm-5mm-2.9cm\relax}
    {\color{cursfosc}\small\bfseries\MakeUppercase{#3}}\par\vspace{2pt}
    {\color{navy}\fontsize{20}{23}\selectfont\bfseries #4\par}
  \end{minipage}\hfill
  \begin{minipage}[c]{2.6cm}\raggedleft\includegraphics[width=2.6cm]{\fm@logo}\end{minipage}\par
  \vspace{2pt}}

% ---------- barres de secció del color del curs ----------
\renewcommand{\seccio}[2]{\par\needspace{\fmneed\baselineskip}\bigskip
  \begin{tcolorbox}[colback=curs,colframe=curs,coltext=curstext,arc=2mm,boxrule=0pt,left=3mm,right=3mm,top=1.2mm,bottom=1.2mm]
  {\large\bfseries #1\quad #2}
  \end{tcolorbox}\vspace{-2pt}}

% ---------- estrelles de dificultat del color del curs (a Batxillerat no se'n posen) ----------
\@ifundefined{estrelles}{}{%
  \renewcommand{\estrelles}[1]{{\color{curs}\ifcase#1\or\faStar\faStar[regular]\faStar[regular]\or\faStar\faStar\faStar[regular]\or\faStar\faStar\faStar\fi}}}

% ---------- capçalera i peu de les pàgines ----------
% primera pàgina: logos institucionals al peu
\fancypagestyle{primera}{\fancyhf{}\renewcommand{\headrulewidth}{0pt}
  \fancyfoot[L]{\raisebox{-0.32\height}{\includegraphics[height=0.78cm]{logo-gva.png}}\hspace{5mm}%
    \raisebox{-0.32\height}{\includegraphics[height=0.78cm]{logo-vermellar.png}}}
  \fancyfoot[R]{\scriptsize\color{muted}\href{https://creativecommons.org/licenses/by-nc-sa/4.0/deed.ca}{\ccbyncsa}\ \ CC BY-NC-SA 4.0 · fentmates.org\qquad\small\thepage}}
% la resta: el símbol a l'esquerra i el curs a la dreta, amb el títol al centre com abans
\fancyhead[L]{\raisebox{-0.3\height}{\includegraphics[height=0.62cm]{\fm@marca}}}
\fancyhead[R]{\small\bfseries\color{cursfosc}\fm@cursnom}
\renewcommand{\headrule}{\color{curs!45}\hrule height \headrulewidth}
\makeatother

%   \fmcurs{1eso}
% i, al començament del document:
%   \capcalera{Matemàtiques}{1r ESO}{Fitxa d'exercicis · Unitat 2}{Divisibilitat}
% Opcional abans de \input: \def\fmneed{14} (espai mínim abans de cada secció).
% =====================================================================
\makeatletter
\def\fm@logodir{../../logo/pdf}
\providecommand{\fmneed}{10}
\def\fm@cursnom{}

% clau del curs, color i color del text damunt del color
\newcommand{\fm@curs}[3]{\@namedef{fm@color@#1}{#2}\@namedef{fm@text@#1}{#3}}
\fm@curs{1eso}{E6B84A}{navy}
\fm@curs{2eso}{3B82F6}{white}
\fm@curs{3eso}{FF7A50}{white}
\fm@curs{4eso-a}{A78BFA}{white}
\fm@curs{4eso-b}{6D28D9}{white}
\fm@curs{1batx-mi}{9B1C1C}{white}
\fm@curs{1batx-ccss}{EF6B6B}{white}
\fm@curs{2batx-mii}{1E6B47}{white}
\fm@curs{2batx-ccss}{2BB3A3}{white}

\newcommand{\fmcurs}[1]{%
  \@ifundefined{fm@color@#1}{\PackageError{fentmates}{Curs desconegut: #1}{}}{}%
  \edef\fm@tmp{\noexpand\definecolor{curs}{HTML}{\@nameuse{fm@color@#1}}%
               \noexpand\colorlet{curstext}{\@nameuse{fm@text@#1}}}\fm@tmp
  \colorlet{cursfosc}{curs!72!black}%
  \def\fm@logo{\fm@logodir/vertical-#1.pdf}%
  \def\fm@marca{\fm@logodir/marca-#1.pdf}}

% ---------- capçalera de la primera pàgina ----------
% #1 matèria (petit), #2 curs, #3 tipus de document i unitat, #4 títol
\newcommand{\capcalera}[4]{%
  \gdef\fm@cursnom{#2}%
  \par\noindent
  \begin{tcolorbox}[colback=curs,colframe=curs,boxrule=0pt,arc=2.2mm,width=3.5cm,height=2.15cm,
      valign=center,left=3mm,right=2mm,top=1mm,bottom=1mm,nobeforeafter,box align=center]
    \color{curstext}{\footnotesize #1}\par\vspace{1pt}{\fontsize{19}{21}\selectfont\bfseries #2\par}%
  \end{tcolorbox}\hspace{5mm}%
  \begin{minipage}[c]{\dimexpr\linewidth-3.5cm-5mm-2.9cm\relax}
    {\color{cursfosc}\small\bfseries\MakeUppercase{#3}}\par\vspace{2pt}
    {\color{navy}\fontsize{20}{23}\selectfont\bfseries #4\par}
  \end{minipage}\hfill
  \begin{minipage}[c]{2.6cm}\raggedleft\includegraphics[width=2.6cm]{\fm@logo}\end{minipage}\par
  \vspace{2pt}}

% ---------- barres de secció del color del curs ----------
\renewcommand{\seccio}[2]{\par\needspace{\fmneed\baselineskip}\bigskip
  \begin{tcolorbox}[colback=curs,colframe=curs,coltext=curstext,arc=2mm,boxrule=0pt,left=3mm,right=3mm,top=1.2mm,bottom=1.2mm]
  {\large\bfseries #1\quad #2}
  \end{tcolorbox}\vspace{-2pt}}

% ---------- estrelles de dificultat del color del curs (a Batxillerat no se'n posen) ----------
\@ifundefined{estrelles}{}{%
  \renewcommand{\estrelles}[1]{{\color{curs}\ifcase#1\or\faStar\faStar[regular]\faStar[regular]\or\faStar\faStar\faStar[regular]\or\faStar\faStar\faStar\fi}}}

% ---------- capçalera i peu de les pàgines ----------
% primera pàgina: logos institucionals al peu
\fancypagestyle{primera}{\fancyhf{}\renewcommand{\headrulewidth}{0pt}
  \fancyfoot[L]{\raisebox{-0.32\height}{\includegraphics[height=0.78cm]{logo-gva.png}}\hspace{5mm}%
    \raisebox{-0.32\height}{\includegraphics[height=0.78cm]{logo-vermellar.png}}}
  \fancyfoot[R]{\scriptsize\color{muted}\href{https://creativecommons.org/licenses/by-nc-sa/4.0/deed.ca}{\ccbyncsa}\ \ CC BY-NC-SA 4.0 · fentmates.org\qquad\small\thepage}}
% la resta: el símbol a l'esquerra i el curs a la dreta, amb el títol al centre com abans
\fancyhead[L]{\raisebox{-0.3\height}{\includegraphics[height=0.62cm]{\fm@marca}}}
\fancyhead[R]{\small\bfseries\color{cursfosc}\fm@cursnom}
\renewcommand{\headrule}{\color{curs!45}\hrule height \headrulewidth}
\makeatother

%   \fmcurs{1eso}
% i, al començament del document:
%   \capcalera{Matemàtiques}{1r ESO}{Fitxa d'exercicis · Unitat 2}{Divisibilitat}
% Opcional abans de \input: \def\fmneed{14} (espai mínim abans de cada secció).
% =====================================================================
\makeatletter
\def\fm@logodir{../../logo/pdf}
\providecommand{\fmneed}{10}
\def\fm@cursnom{}

% clau del curs, color i color del text damunt del color
\newcommand{\fm@curs}[3]{\@namedef{fm@color@#1}{#2}\@namedef{fm@text@#1}{#3}}
\fm@curs{1eso}{E6B84A}{navy}
\fm@curs{2eso}{3B82F6}{white}
\fm@curs{3eso}{FF7A50}{white}
\fm@curs{4eso-a}{A78BFA}{white}
\fm@curs{4eso-b}{6D28D9}{white}
\fm@curs{1batx-mi}{9B1C1C}{white}
\fm@curs{1batx-ccss}{EF6B6B}{white}
\fm@curs{2batx-mii}{1E6B47}{white}
\fm@curs{2batx-ccss}{2BB3A3}{white}

\newcommand{\fmcurs}[1]{%
  \@ifundefined{fm@color@#1}{\PackageError{fentmates}{Curs desconegut: #1}{}}{}%
  \edef\fm@tmp{\noexpand\definecolor{curs}{HTML}{\@nameuse{fm@color@#1}}%
               \noexpand\colorlet{curstext}{\@nameuse{fm@text@#1}}}\fm@tmp
  \colorlet{cursfosc}{curs!72!black}%
  \def\fm@logo{\fm@logodir/vertical-#1.pdf}%
  \def\fm@marca{\fm@logodir/marca-#1.pdf}}

% ---------- capçalera de la primera pàgina ----------
% #1 matèria (petit), #2 curs, #3 tipus de document i unitat, #4 títol
\newcommand{\capcalera}[4]{%
  \gdef\fm@cursnom{#2}%
  \par\noindent
  \begin{tcolorbox}[colback=curs,colframe=curs,boxrule=0pt,arc=2.2mm,width=3.5cm,height=2.15cm,
      valign=center,left=3mm,right=2mm,top=1mm,bottom=1mm,nobeforeafter,box align=center]
    \color{curstext}{\footnotesize #1}\par\vspace{1pt}{\fontsize{19}{21}\selectfont\bfseries #2\par}%
  \end{tcolorbox}\hspace{5mm}%
  \begin{minipage}[c]{\dimexpr\linewidth-3.5cm-5mm-2.9cm\relax}
    {\color{cursfosc}\small\bfseries\MakeUppercase{#3}}\par\vspace{2pt}
    {\color{navy}\fontsize{20}{23}\selectfont\bfseries #4\par}
  \end{minipage}\hfill
  \begin{minipage}[c]{2.6cm}\raggedleft\includegraphics[width=2.6cm]{\fm@logo}\end{minipage}\par
  \vspace{2pt}}

% ---------- barres de secció del color del curs ----------
\renewcommand{\seccio}[2]{\par\needspace{\fmneed\baselineskip}\bigskip
  \begin{tcolorbox}[colback=curs,colframe=curs,coltext=curstext,arc=2mm,boxrule=0pt,left=3mm,right=3mm,top=1.2mm,bottom=1.2mm]
  {\large\bfseries #1\quad #2}
  \end{tcolorbox}\vspace{-2pt}}

% ---------- estrelles de dificultat del color del curs (a Batxillerat no se'n posen) ----------
\@ifundefined{estrelles}{}{%
  \renewcommand{\estrelles}[1]{{\color{curs}\ifcase#1\or\faStar\faStar[regular]\faStar[regular]\or\faStar\faStar\faStar[regular]\or\faStar\faStar\faStar\fi}}}

% ---------- capçalera i peu de les pàgines ----------
% primera pàgina: logos institucionals al peu
\fancypagestyle{primera}{\fancyhf{}\renewcommand{\headrulewidth}{0pt}
  \fancyfoot[L]{\raisebox{-0.32\height}{\includegraphics[height=0.78cm]{logo-gva.png}}\hspace{5mm}%
    \raisebox{-0.32\height}{\includegraphics[height=0.78cm]{logo-vermellar.png}}}
  \fancyfoot[R]{\scriptsize\color{muted}\href{https://creativecommons.org/licenses/by-nc-sa/4.0/deed.ca}{\ccbyncsa}\ \ CC BY-NC-SA 4.0 · fentmates.org\qquad\small\thepage}}
% la resta: el símbol a l'esquerra i el curs a la dreta, amb el títol al centre com abans
\fancyhead[L]{\raisebox{-0.3\height}{\includegraphics[height=0.62cm]{\fm@marca}}}
\fancyhead[R]{\small\bfseries\color{cursfosc}\fm@cursnom}
\renewcommand{\headrule}{\color{curs!45}\hrule height \headrulewidth}
\makeatother

\fmcurs{4eso-a}

\begin{document}
\thispagestyle{primera}

\capcalera{Matemáticas A}{4\textsuperscript{o} ESO}{Soluciones de la ficha · Unidad 11}{Probabilidad}
\vspace{8pt}

\begin{tcolorbox}[colback=soft,colframe=soft,arc=2mm,left=3mm,right=3mm,top=1.5mm,bottom=1.5mm]
\small Intenta resolver cada ejercicio antes de mirar su solución. Si no te sale el mismo resultado, repasa los pasos o consulta los apuntes interactivos de la unidad: encontrarás simulaciones, una bolsa de bolas, la tabla de los dos dados, diagramas de árbol de experiencias independientes y dependientes, dos urnas y una tabla de contingencia que calculan las probabilidades paso a paso. Las probabilidades se dan en fracción irreducible (o en decimal, si los datos lo son).
\end{tcolorbox}

% =====================================================
\seccio{1}{Experimentos}

\ex{1} \begin{tasks}(2)
  \task \rest{aleatorio}
  \task \rest{determinista} (100 °C)
  \task \rest{aleatorio}
  \task \rest{determinista} (6 €)
  \task \rest{aleatorio}
  \task \rest{determinista}
\end{tasks}

\ex{1} \begin{tasks}(4)
  \task \rest{muy probable}
  \task \rest{poco probable}
  \task \rest{imposible}
  \task \rest{seguro}
\end{tasks}

\ex{1} \begin{tasks}(2)
  \task $E=\{C,\,X\}$
  \task $E=\{1,\,2,\,3,\,4\}$
  \task $E=\{\text{rojo},\,\text{azul},\,\text{verde},\,\text{amarillo},\,\text{negro}\}$
  \task $E=\{R,\,A\}$
  \task $E=\{\text{L},\,\text{M},\,\text{X},\,\text{J},\,\text{V},\,\text{S},\,\text{D}\}$
  \task $E=\{0,\,1,\,2,\,3\}$
\end{tasks}

\ex{2} \rest{No}. Sabemos qué resultados \emph{pueden} salir (el espacio muestral), pero no sabemos \emph{cuál} saldrá en cada lanzamiento: es un experimento \textbf{aleatorio}.

\ex{2} \par\smallskip\begin{center}\begin{tikzpicture}[grow'=right,level distance=14mm,level 1/.style={sibling distance=13mm},level 2/.style={sibling distance=6.5mm},every node/.style={font=\small\bfseries},edge from parent/.style={draw=muted,thick}]
\node[circle,fill=navy,inner sep=1.2pt] {} child {node[] {C} child {node[label={[font=\scriptsize,xshift=1mm]right:$CC$}] {C} edge from parent node[above,pos=.55,font=\scriptsize,text=teal!70!black,inner sep=1.5pt] {}} child {node[label={[font=\scriptsize,xshift=1mm]right:$CX$}] {X} edge from parent node[below,pos=.55,font=\scriptsize,text=teal!70!black,inner sep=1.5pt] {}} edge from parent node[above,pos=.55,font=\scriptsize,text=teal!70!black,inner sep=1.5pt] {}} child {node[] {X} child {node[label={[font=\scriptsize,xshift=1mm]right:$XC$}] {C} edge from parent node[above,pos=.55,font=\scriptsize,text=teal!70!black,inner sep=1.5pt] {}} child {node[label={[font=\scriptsize,xshift=1mm]right:$XX$}] {X} edge from parent node[below,pos=.55,font=\scriptsize,text=teal!70!black,inner sep=1.5pt] {}} edge from parent node[below,pos=.55,font=\scriptsize,text=teal!70!black,inner sep=1.5pt] {}};
\end{tikzpicture}\end{center}$E=\{CC,\,CX,\,XC,\,XX\}$; \quad $A=\res{\{CC,\,CX,\,XC\}}$, \quad $B=\res{\{XX\}}$, \quad $C=\res{\{CC,\,XX\}}$.

\ex{2} \par\smallskip\begin{center}\begin{tikzpicture}[grow'=right,level distance=13mm,level 1/.style={sibling distance=20mm},level 2/.style={sibling distance=10mm},level 3/.style={sibling distance=5mm},every node/.style={font=\small\bfseries},edge from parent/.style={draw=muted,thick}]
\node[circle,fill=navy,inner sep=1.2pt] {} child {node[] {O} child {node[] {O} child {node[label={[font=\scriptsize,xshift=1mm]right:$OOO$}] {O} edge from parent node[above,pos=.55,font=\scriptsize,text=teal!70!black,inner sep=1.5pt] {}} child {node[label={[font=\scriptsize,xshift=1mm]right:$OOA$}] {A} edge from parent node[below,pos=.55,font=\scriptsize,text=teal!70!black,inner sep=1.5pt] {}} edge from parent node[above,pos=.55,font=\scriptsize,text=teal!70!black,inner sep=1.5pt] {}} child {node[] {A} child {node[label={[font=\scriptsize,xshift=1mm]right:$OAO$}] {O} edge from parent node[above,pos=.55,font=\scriptsize,text=teal!70!black,inner sep=1.5pt] {}} child {node[label={[font=\scriptsize,xshift=1mm]right:$OAA$}] {A} edge from parent node[below,pos=.55,font=\scriptsize,text=teal!70!black,inner sep=1.5pt] {}} edge from parent node[below,pos=.55,font=\scriptsize,text=teal!70!black,inner sep=1.5pt] {}} edge from parent node[above,pos=.55,font=\scriptsize,text=teal!70!black,inner sep=1.5pt] {}} child {node[] {A} child {node[] {O} child {node[label={[font=\scriptsize,xshift=1mm]right:$AOO$}] {O} edge from parent node[above,pos=.55,font=\scriptsize,text=teal!70!black,inner sep=1.5pt] {}} child {node[label={[font=\scriptsize,xshift=1mm]right:$AOA$}] {A} edge from parent node[below,pos=.55,font=\scriptsize,text=teal!70!black,inner sep=1.5pt] {}} edge from parent node[above,pos=.55,font=\scriptsize,text=teal!70!black,inner sep=1.5pt] {}} child {node[] {A} child {node[label={[font=\scriptsize,xshift=1mm]right:$AAO$}] {O} edge from parent node[above,pos=.55,font=\scriptsize,text=teal!70!black,inner sep=1.5pt] {}} child {node[label={[font=\scriptsize,xshift=1mm]right:$AAA$}] {A} edge from parent node[below,pos=.55,font=\scriptsize,text=teal!70!black,inner sep=1.5pt] {}} edge from parent node[below,pos=.55,font=\scriptsize,text=teal!70!black,inner sep=1.5pt] {}} edge from parent node[below,pos=.55,font=\scriptsize,text=teal!70!black,inner sep=1.5pt] {}};
\end{tikzpicture}\end{center}$E=\{OOO,\,OOA,\,OAO,\,OAA,\,AOO,\,AOA,\,AAO,\,AAA\}$ (8 resultados).\par «Dos niñas» $=\res{\{OAA,\,AOA,\,AAO\}}$.

\ex{2} \par\smallskip\begin{center}\begin{tikzpicture}[grow'=right,level distance=14mm,level 1/.style={sibling distance=15mm},level 2/.style={sibling distance=6mm},every node/.style={font=\small\bfseries},edge from parent/.style={draw=muted,thick}]
\node[circle,fill=navy,inner sep=1.2pt] {} child {node[] {S} child {node[label={[font=\scriptsize,xshift=1mm]right:$SC$}] {C} edge from parent node[above,pos=.55,font=\scriptsize,text=teal!70!black,inner sep=1.5pt] {}} child {node[label={[font=\scriptsize,xshift=1mm]right:$SP$}] {P} edge from parent node[above,pos=.55,font=\scriptsize,text=teal!70!black,inner sep=1.5pt] {}} child {node[label={[font=\scriptsize,xshift=1mm]right:$ST$}] {T} edge from parent node[below,pos=.55,font=\scriptsize,text=teal!70!black,inner sep=1.5pt] {}} edge from parent node[above,pos=.55,font=\scriptsize,text=teal!70!black,inner sep=1.5pt] {}} child {node[] {E} child {node[label={[font=\scriptsize,xshift=1mm]right:$EC$}] {C} edge from parent node[above,pos=.55,font=\scriptsize,text=teal!70!black,inner sep=1.5pt] {}} child {node[label={[font=\scriptsize,xshift=1mm]right:$EP$}] {P} edge from parent node[above,pos=.55,font=\scriptsize,text=teal!70!black,inner sep=1.5pt] {}} child {node[label={[font=\scriptsize,xshift=1mm]right:$ET$}] {T} edge from parent node[below,pos=.55,font=\scriptsize,text=teal!70!black,inner sep=1.5pt] {}} edge from parent node[below,pos=.55,font=\scriptsize,text=teal!70!black,inner sep=1.5pt] {}};
\end{tikzpicture}\end{center}$E=\{SC,\,SP,\,ST,\,EC,\,EP,\,ET\}$ ($2\cdot3=6$ menús). \quad «Pescado» $=\res{\{SP,\,EP\}}$.

\ex{3} $E=\{123,\,132,\,213,\,231,\,312,\,321\}$.\par «Par» $=\res{\{132,\,312\}}$; \quad «mayor que 200» $=\res{\{213,\,231,\,312,\,321\}}$.


% =====================================================
\seccio{2}{Sucesos}

\ex{1} \begin{tasks}(3)
  \task $\{2,4,6,8,10,12\}$: \rest{compuesto}
  \task $\{5,\,10\}$: \rest{compuesto}
  \task $\varnothing$: \rest{imposible}
  \task $\{1,2,3,4,6,12\}$: \rest{compuesto}
  \task $E$: \rest{seguro}
  \task $\{11\}$: \rest{elemental}
\end{tasks}

\ex{1} $E=\{1,\,2,\,3,\,4,\,5,\,6,\,7,\,8\}$\begin{tasks}(4)
  \task $\res{\{2,4,6,8\}}$
  \task $\res{\{2,3,5,7\}}$
  \task $\res{\{6,7,8\}}$
  \task $\res{\{1,2,4,8\}}$
\end{tasks}

\ex{1} $A=\{3,\,6\}$, $B=\{1,\,2\}$: \ $A\cup B=\res{\{1,\,2,\,3,\,6\}}$, \ $A\cap B=\res{\varnothing}$, \ $\overline{A}=\res{\{1,\,2,\,4,\,5\}}$, \ $\overline{B}=\res{\{3,\,4,\,5,\,6\}}$. \ Son \rest{incompatibles}, pero \rest{no son contrarios}: si sale un 4 o un 5, no ocurre ninguno de los dos.

\ex{1} $A=\{2,4,6,8,10\}$, $B=\{2,3,5,7\}$: \ $A\cup B=\res{\{2,3,4,5,6,7,8,10\}}$, \ $A\cap B=\res{\{2\}}$, \ $\overline{A}=\res{\{1,3,5,7,9\}}$, \ $\overline{B}=\res{\{1,4,6,8,9,10\}}$.

\ex{2} $O\cap A$ = el as de oros: \res{1 carta}. \ $O\cup A$ = los 10 oros y los otros 3 ases: $10+4-1=\res{13\text{ cartas}}$. \ $\overline{O}$ = las que no son de oros: $40-10=\res{30\text{ cartas}}$.

\ex{2} $A=\{CCC,\,CCX,\,CXC,\,XCC\}$, \ $B=\{XCC,\,XCX,\,XXC,\,XXX\}$. \ $A\cap B=\res{\{XCC\}}$; \ $A\cup B$: todos los resultados menos $CXX$ (\res{7 resultados}); \ $\overline{A}=\res{\{CXX,\,XCX,\,XXC,\,XXX\}}$.

\ex{3} $A=\{2,4,6,8,10,12\}$, $B=\{3,6,9,12\}$.
\begin{tasks}(1)
  \task $A\cap\overline{B}=\res{\{2,4,8,10\}}$.
  \task $\overline{A}\cap\overline{B}=\res{\{1,5,7,11\}}$. \ $A\cup B=\{2,3,4,6,8,9,10,12\}$, y su contrario es $\{1,5,7,11\}$ \checkmark
\end{tasks}


% =====================================================
\seccio{3}{Probabilidad de un suceso}

\ex{1} $h$: \ $\res{0{,}14}$, \ $\res{0{,}18}$, \ $\res{0{,}16}$, \ $\res{0{,}2}$, \ $\res{0{,}12}$, \ $\res{0{,}2}$ \ (suman 1). \quad $h(\text{par})=\frac{9+10+10}{50}=\res{0{,}58}$.

\ex{1} Una probabilidad siempre está entre 0 y 1. Pueden serlo: \rest{$0{,}7$, \ $0$, \ $100\,\%$ y $\frac{99}{100}$}. No pueden serlo: $1{,}2$, $-0{,}1$, $\frac32$ y $\frac54$.

\ex{1} \begin{tasks}(3)
  \task $P(\overline{A})=\res{0{,}38}$
  \task $P(\overline{B})=\res{\frac47}$
  \task $P(\overline{C})=\res{85\,\%}$
\end{tasks}

\ex{2} $h$: \ $0{,}65$; \ $0{,}58$; \ $0{,}594$; \ $0{,}601$. \ Con más lanzamientos se estabiliza: $P(\text{punta arriba})\approx\res{0{,}6}$.

\ex{2} $P\approx\frac{45}{60}=\res{0{,}75}$. \ $20\cdot0{,}75=\res{15\text{ penaltis}}$ (aproximadamente).

\ex{2} $P(6)=1-(0{,}15+0{,}1+0{,}2+0{,}15+0{,}1)=1-0{,}7=\res{0{,}3}$. \ $P(\text{par})=P(2)+P(4)+P(6)=0{,}1+0{,}15+0{,}3=\res{0{,}55}$.

\ex{2} $P(A\cup B)=0{,}4+0{,}5-0{,}15=\res{0{,}75}$. \ $P(\overline{A\cup B})=1-0{,}75=\res{0{,}25}$. \ $P(A\cap\overline{B})=0{,}4-0{,}15=\res{0{,}25}$.

\ex{2} 
\begin{tasks}(1)
  \task $P(A\cup B)=0{,}3+0{,}45=\res{0{,}75}$.
  \task \rest{No}: si fueran incompatibles, $P(A\cup B)=0{,}6+0{,}5=1{,}1>1$, y eso es imposible.
\end{tasks}

\ex{3} $P(R)\approx\frac{120}{300}=0{,}4$: más o menos, el 40 \% de las bolas son rojas. \ $30\cdot0{,}4=\res{12\text{ rojas}}$ (y 18 blancas).

\ex{3} $P(A)=1-0{,}3=\res{0{,}7}$; \ $P(B)=1-0{,}4=\res{0{,}6}$. \ De la fórmula de la unión: $0{,}9=0{,}7+0{,}6-P(A\cap B)$, \ $P(A\cap B)=\res{0{,}4}$. \ Como $P(A\cap B)\neq0$, \rest{no son incompatibles}.


% =====================================================
\seccio{4}{Cálculo de probabilidades}

\ex{1} \begin{tasks}(3)
  \task $\res{\frac16}$
  \task $\frac36=\res{\frac12}$
  \task $\frac46=\res{\frac23}$
  \task $\frac26=\res{\frac13}$
  \task $\res{0}$ (imposible)
  \task $\res{1}$ (seguro)
\end{tasks}

\ex{1} Hay 15 bolas: \ $P(\text{roja})=\frac6{15}=\res{\frac25}$, \ $P(\text{azul})=\res{\frac4{15}}$, \ $P(\text{verde})=\frac5{15}=\res{\frac13}$, \ $P(\text{no verde})=\frac{6+4}{15}=\res{\frac23}$.

\ex{1} \par\smallskip\begin{center}{\small\renewcommand{\arraystretch}{1.15}\setlength{\tabcolsep}{4pt}\begin{tabular}{c|cccccc}
\rowcolor{soft} $+$ & 1 & 2 & 3 & 4 & 5 & 6\\\hline
\cellcolor{soft}1 & 2 & 3 & 4 & 5 & 6 & 7\\
\cellcolor{soft}2 & 3 & 4 & 5 & 6 & 7 & \cellcolor{accent!30}\textbf{8}\\
\cellcolor{soft}3 & 4 & 5 & 6 & 7 & \cellcolor{accent!30}\textbf{8} & 9\\
\cellcolor{soft}4 & 5 & 6 & 7 & \cellcolor{accent!30}\textbf{8} & 9 & 10\\
\cellcolor{soft}5 & 6 & 7 & \cellcolor{accent!30}\textbf{8} & 9 & 10 & 11\\
\cellcolor{soft}6 & 7 & \cellcolor{accent!30}\textbf{8} & 9 & 10 & 11 & 12\\
\end{tabular}}\end{center}
\begin{tasks}(1)
  \task $\res{\frac5{36}}$ (marcadas).
  \task $\res{\frac1{36}}$.
  \task $\frac{3+2+1}{36}=\res{\frac16}$.
  \task $\frac{18}{36}=\res{\frac12}$.
\end{tasks}

\ex{2} \begin{tasks}(3)
  \task $\frac{10}{40}=\res{\frac14}$
  \task $\frac4{40}=\res{\frac1{10}}$
  \task $\res{\frac3{40}}$
  \task $\frac{28}{40}=\res{\frac7{10}}$
  \task $\frac{10+3}{40}=\res{\frac{13}{40}}$ (el as de oros, una vez)
  \task $\res{\frac1{40}}$
\end{tasks}

\ex{2} 11 letras; vocales A, E, A, I, A: \ $P(\text{vocal})=\res{\frac{5}{11}}$; \ $P(M)=\res{\frac{2}{11}}$; \ $P(\text{consonante})=\res{\frac{6}{11}}$.

\ex{2} $P(\text{Pablo})=\frac26=\frac13$ y $P(\text{María})=\frac46=\frac23$: \rest{no es justo}. Por ejemplo: Pablo gana con par y María con impar ($\frac12$ cada uno).

\ex{2} 
\begin{tasks}(1)
  \task $\frac6{36}=\res{\frac16}$.
  \task $(1,3),(2,4),(3,5),(4,6)$ y al revés: $\frac8{36}=\res{\frac29}$.
  \task Con la tabla de los productos, 27 casillas son pares: $\frac{27}{36}=\res{\frac34}$.
\end{tasks}

\ex{3} Si $P(\text{azul})=\frac23$, la roja tiene probabilidad $\frac13$: 1 de cada 3 bolas es roja, así que hay $5\cdot3=15$ bolas y $15-5=\res{10\text{ azules}}$. \ (Comprobación: $\frac{10}{15}=\frac23$ \checkmark)

\ex{3} \rest{No} es regular: los sectores no son iguales. La probabilidad es la fracción de círculo que ocupa cada sector: \ $P(\text{rojo})=\frac{180}{360}=\res{\frac12}$, \ $P(\text{azul})=\frac{120}{360}=\res{\frac13}$, \ $P(\text{verde})=\frac{60}{360}=\res{\frac16}$.


% =====================================================
\seccio{5}{Experiencias compuestas. Probabilidad}

\ex{1} \par\smallskip\begin{center}\begin{tikzpicture}[grow'=right,level distance=13mm,level 1/.style={sibling distance=24mm},level 2/.style={sibling distance=3.7mm},every node/.style={font=\scriptsize\bfseries},edge from parent/.style={draw=muted,thick}]
\node[circle,fill=navy,inner sep=1.2pt] {} child {node[] {C} child {node[label={[font=\scriptsize,xshift=1mm]right:$C1$}] {1} edge from parent node[above,pos=.55,font=\scriptsize,text=teal!70!black,inner sep=1.5pt] {}} child {node[label={[font=\scriptsize,xshift=1mm]right:$C2$}] {2} edge from parent node[above,pos=.55,font=\scriptsize,text=teal!70!black,inner sep=1.5pt] {}} child {node[label={[font=\scriptsize,xshift=1mm]right:$C3$}] {3} edge from parent node[above,pos=.55,font=\scriptsize,text=teal!70!black,inner sep=1.5pt] {}} child {node[label={[font=\scriptsize,xshift=1mm]right:$C4$}] {4} edge from parent node[below,pos=.55,font=\scriptsize,text=teal!70!black,inner sep=1.5pt] {}} child {node[label={[font=\scriptsize,xshift=1mm]right:$C5$}] {5} edge from parent node[below,pos=.55,font=\scriptsize,text=teal!70!black,inner sep=1.5pt] {}} child {node[label={[font=\scriptsize,xshift=1mm]right:$C6$}] {6} edge from parent node[below,pos=.55,font=\scriptsize,text=teal!70!black,inner sep=1.5pt] {}} edge from parent node[above,pos=.55,font=\scriptsize,text=teal!70!black,inner sep=1.5pt] {}} child {node[] {X} child {node[label={[font=\scriptsize,xshift=1mm]right:$X1$}] {1} edge from parent node[above,pos=.55,font=\scriptsize,text=teal!70!black,inner sep=1.5pt] {}} child {node[label={[font=\scriptsize,xshift=1mm]right:$X2$}] {2} edge from parent node[above,pos=.55,font=\scriptsize,text=teal!70!black,inner sep=1.5pt] {}} child {node[label={[font=\scriptsize,xshift=1mm]right:$X3$}] {3} edge from parent node[above,pos=.55,font=\scriptsize,text=teal!70!black,inner sep=1.5pt] {}} child {node[label={[font=\scriptsize,xshift=1mm]right:$X4$}] {4} edge from parent node[below,pos=.55,font=\scriptsize,text=teal!70!black,inner sep=1.5pt] {}} child {node[label={[font=\scriptsize,xshift=1mm]right:$X5$}] {5} edge from parent node[below,pos=.55,font=\scriptsize,text=teal!70!black,inner sep=1.5pt] {}} child {node[label={[font=\scriptsize,xshift=1mm]right:$X6$}] {6} edge from parent node[below,pos=.55,font=\scriptsize,text=teal!70!black,inner sep=1.5pt] {}} edge from parent node[below,pos=.55,font=\scriptsize,text=teal!70!black,inner sep=1.5pt] {}};
\end{tikzpicture}\end{center}$2\cdot6=\res{12}$ resultados equiprobables.
\begin{tasks}(1)
  \task $\frac1{12}$ (o bien $\frac12\cdot\frac16$) $=\res{\frac1{12}}$.
  \task $X2,X4,X6$: $\frac3{12}=\res{\frac14}$.
  \task $C2,C3,C5,X2,X3,X5$: $\frac6{12}=\res{\frac12}$.
\end{tasks}

\ex{1} \par\smallskip\begin{center}\begin{tikzpicture}[grow'=right,level distance=16mm,level 1/.style={sibling distance=14mm},level 2/.style={sibling distance=7mm},every node/.style={font=\small\bfseries},edge from parent/.style={draw=muted,thick}]
\node[circle,fill=navy,inner sep=1.2pt] {} child {node[] {C} child {node[label={[font=\scriptsize,xshift=1mm]right:$CC$: $\frac14$}] {C} edge from parent node[above,pos=.55,font=\scriptsize,text=teal!70!black,inner sep=1.5pt] {$\frac12$}} child {node[label={[font=\scriptsize,xshift=1mm]right:$CX$: $\frac14$}] {X} edge from parent node[below,pos=.55,font=\scriptsize,text=teal!70!black,inner sep=1.5pt] {$\frac12$}} edge from parent node[above,pos=.55,font=\scriptsize,text=teal!70!black,inner sep=1.5pt] {$\frac12$}} child {node[] {X} child {node[label={[font=\scriptsize,xshift=1mm]right:$XC$: $\frac14$}] {C} edge from parent node[above,pos=.55,font=\scriptsize,text=teal!70!black,inner sep=1.5pt] {$\frac12$}} child {node[label={[font=\scriptsize,xshift=1mm]right:$XX$: $\frac14$}] {X} edge from parent node[below,pos=.55,font=\scriptsize,text=teal!70!black,inner sep=1.5pt] {$\frac12$}} edge from parent node[below,pos=.55,font=\scriptsize,text=teal!70!black,inner sep=1.5pt] {$\frac12$}};
\end{tikzpicture}\end{center}$P(CC)=\frac12\cdot\frac12=\res{\frac14}$. \ «Al menos una cruz» son 3 caminos ($CX$, $XC$, $XX$): $\frac14+\frac14+\frac14=\res{\frac34}$.

\ex{1} \par\smallskip\begin{center}\begin{tikzpicture}[grow'=right,level distance=17mm,level 1/.style={sibling distance=15mm},level 2/.style={sibling distance=7mm},every node/.style={font=\small\bfseries},edge from parent/.style={draw=muted,thick}]
\node[circle,fill=navy,inner sep=1.2pt] {} child {node[] {D} child {node[label={[font=\scriptsize,xshift=1mm]right:$0{,}18$}] {D} edge from parent node[above,pos=.55,font=\scriptsize,text=teal!70!black,inner sep=1.5pt] {$0{,}6$}} child {node[label={[font=\scriptsize,xshift=1mm]right:$0{,}12$}] {$\overline{D}$} edge from parent node[below,pos=.55,font=\scriptsize,text=teal!70!black,inner sep=1.5pt] {$0{,}4$}} edge from parent node[above,pos=.55,font=\scriptsize,text=teal!70!black,inner sep=1.5pt] {$0{,}3$}} child {node[] {$\overline{D}$} child {node[label={[font=\scriptsize,xshift=1mm]right:$0{,}42$}] {D} edge from parent node[above,pos=.55,font=\scriptsize,text=teal!70!black,inner sep=1.5pt] {$0{,}6$}} child {node[label={[font=\scriptsize,xshift=1mm]right:$0{,}28$}] {$\overline{D}$} edge from parent node[below,pos=.55,font=\scriptsize,text=teal!70!black,inner sep=1.5pt] {$0{,}4$}} edge from parent node[below,pos=.55,font=\scriptsize,text=teal!70!black,inner sep=1.5pt] {$0{,}7$}};
\end{tikzpicture}\end{center}
\begin{tasks}(1)
  \task $0{,}3\cdot0{,}6=\res{0{,}18}$.
  \task $0{,}7\cdot0{,}4=\res{0{,}28}$.
  \task $1-0{,}28=\res{0{,}72}$.
  \task $0{,}3\cdot0{,}4+0{,}7\cdot0{,}6=0{,}12+0{,}42=\res{0{,}54}$.
\end{tasks}

\ex{2} Cada camino del árbol tiene probabilidad $\frac12\cdot\frac12\cdot\frac12=\frac18$.
\begin{tasks}(1)
  \task $\res{\frac18}$.
  \task $CXX$, $XCX$, $XXC$: $\res{\frac38}$.
  \task $CCC$, $CCX$, $CXC$, $XCC$: $\frac48=\res{\frac12}$.
\end{tasks}

\ex{2} \par\smallskip\begin{center}\begin{tikzpicture}[grow'=right,level distance=17mm,level 1/.style={sibling distance=15mm},level 2/.style={sibling distance=7.5mm},every node/.style={font=\small\bfseries},edge from parent/.style={draw=muted,thick}]
\node[circle,fill=navy,inner sep=1.2pt] {} child {node[] {G} child {node[label={[font=\scriptsize,xshift=1mm]right:$GG$: $0{,}64$}] {G} edge from parent node[above,pos=.55,font=\scriptsize,text=teal!70!black,inner sep=1.5pt] {$0{,}8$}} child {node[label={[font=\scriptsize,xshift=1mm]right:$GF$: $0{,}16$}] {F} edge from parent node[below,pos=.55,font=\scriptsize,text=teal!70!black,inner sep=1.5pt] {$0{,}2$}} edge from parent node[above,pos=.55,font=\scriptsize,text=teal!70!black,inner sep=1.5pt] {$0{,}8$}} child {node[] {F} child {node[label={[font=\scriptsize,xshift=1mm]right:$FG$: $0{,}16$}] {G} edge from parent node[above,pos=.55,font=\scriptsize,text=teal!70!black,inner sep=1.5pt] {$0{,}8$}} child {node[label={[font=\scriptsize,xshift=1mm]right:$FF$: $0{,}04$}] {F} edge from parent node[below,pos=.55,font=\scriptsize,text=teal!70!black,inner sep=1.5pt] {$0{,}2$}} edge from parent node[below,pos=.55,font=\scriptsize,text=teal!70!black,inner sep=1.5pt] {$0{,}2$}};
\end{tikzpicture}\end{center}$P(GG)=\res{0{,}64}$; \ $P(FF)=\res{0{,}04}$; \ $P(\text{exactamente uno})=0{,}16+0{,}16=\res{0{,}32}$. \ (Suman 1 \checkmark)

\ex{2} \par\smallskip\begin{center}\begin{tikzpicture}[grow'=right,level distance=17mm,level 1/.style={sibling distance=16mm},level 2/.style={sibling distance=8mm},every node/.style={font=\small\bfseries},edge from parent/.style={draw=muted,thick}]
\node[circle,fill=navy,inner sep=1.2pt] {} child {node[circle,fill=blue!70!black,text=white,inner sep=1.2pt,minimum size=4.2mm,font=\scriptsize\bfseries] {B} child {node[circle,fill=blue!70!black,text=white,inner sep=1.2pt,minimum size=4.2mm,font=\scriptsize\bfseries,label={[font=\scriptsize,xshift=1mm]right:$BB$: $\frac{9}{25}$}] {B} edge from parent node[above,pos=.55,font=\scriptsize,text=teal!70!black,inner sep=1.5pt] {$\frac{3}{5}$}} child {node[circle,fill=black!80,text=white,inner sep=1.2pt,minimum size=4.2mm,font=\scriptsize\bfseries,label={[font=\scriptsize,xshift=1mm]right:$BN$: $\frac{6}{25}$}] {N} edge from parent node[below,pos=.55,font=\scriptsize,text=teal!70!black,inner sep=1.5pt] {$\frac{2}{5}$}} edge from parent node[above,pos=.55,font=\scriptsize,text=teal!70!black,inner sep=1.5pt] {$\frac{3}{5}$}} child {node[circle,fill=black!80,text=white,inner sep=1.2pt,minimum size=4.2mm,font=\scriptsize\bfseries] {N} child {node[circle,fill=blue!70!black,text=white,inner sep=1.2pt,minimum size=4.2mm,font=\scriptsize\bfseries,label={[font=\scriptsize,xshift=1mm]right:$NB$: $\frac{6}{25}$}] {B} edge from parent node[above,pos=.55,font=\scriptsize,text=teal!70!black,inner sep=1.5pt] {$\frac{3}{5}$}} child {node[circle,fill=black!80,text=white,inner sep=1.2pt,minimum size=4.2mm,font=\scriptsize\bfseries,label={[font=\scriptsize,xshift=1mm]right:$NN$: $\frac{4}{25}$}] {N} edge from parent node[below,pos=.55,font=\scriptsize,text=teal!70!black,inner sep=1.5pt] {$\frac{2}{5}$}} edge from parent node[below,pos=.55,font=\scriptsize,text=teal!70!black,inner sep=1.5pt] {$\frac{2}{5}$}};
\end{tikzpicture}\end{center}
\begin{tasks}(1)
  \task $\frac35\cdot\frac35=\res{\frac9{25}}$.
  \task $\frac25\cdot\frac25=\res{\frac4{25}}$.
  \task $\frac6{25}+\frac6{25}=\res{\frac{12}{25}}$.
\end{tasks}

\ex{2} $P(\text{ningún 6})=\frac56\cdot\frac56\cdot\frac56=\res{\frac{125}{216}}$. \ $P(\text{al menos un 6})=1-\frac{125}{216}=\res{\frac{91}{216}}\approx0{,}42$. \ (No es $\frac16+\frac16+\frac16$: esa suma cuenta más de una vez los casos con dos o tres seises.)

\ex{3} En el árbol, cada pregunta tiene las ramas A ($\frac14$) y F ($\frac34$). \ $P(AAA)=\frac14\cdot\frac14\cdot\frac14=\res{\frac1{64}}$. \ $P(FFF)=\frac34\cdot\frac34\cdot\frac34=\res{\frac{27}{64}}$. \ Todos los caminos suman 1: $P(\text{al menos una})=1-\frac{27}{64}=\res{\frac{37}{64}}$.

\ex{3} Con $n$ lanzamientos, $P(\text{al menos un 6})=1-\left(\frac56\right)^n$. \ $n=3$: $1-\frac{125}{216}\approx0{,}42$; \ $n=4$: $1-\frac{625}{1296}=\frac{671}{1296}\approx0{,}52$. \ Hay que lanzarlo \res{4\text{ veces}}.


% =====================================================
\seccio{6}{Experiencias dependientes. Probabilidad condicionada}

\ex{1} Quedan 7 bolas: 4 rojas y 3 verdes. \ $P(R_2/R_1)=\res{\frac47}$; \ $P(V_2/R_1)=\res{\frac37}$.

\ex{1} 
\begin{tasks}(1)
  \task Quedan 9 oros entre 39 cartas: $P(O_2/O_1)=\frac9{39}=\res{\frac3{13}}$.
  \task $P(O_1\cap O_2)=\frac{10}{40}\cdot\frac9{39}=\frac{90}{1560}=\res{\frac3{52}}$.
\end{tasks}

\ex{2} \par\smallskip\begin{center}\begin{tikzpicture}[grow'=right,level distance=17mm,level 1/.style={sibling distance=16mm},level 2/.style={sibling distance=8mm},every node/.style={font=\small\bfseries},edge from parent/.style={draw=muted,thick}]
\node[circle,fill=navy,inner sep=1.2pt] {} child {node[circle,fill=blue!70!black,text=white,inner sep=1.2pt,minimum size=4.2mm,font=\scriptsize\bfseries] {B} child {node[circle,fill=blue!70!black,text=white,inner sep=1.2pt,minimum size=4.2mm,font=\scriptsize\bfseries,label={[font=\scriptsize,xshift=1mm]right:$BB$: $\frac{3}{10}$}] {B} edge from parent node[above,pos=.55,font=\scriptsize,text=teal!70!black,inner sep=1.5pt] {$\frac{2}{4}$}} child {node[circle,fill=black!80,text=white,inner sep=1.2pt,minimum size=4.2mm,font=\scriptsize\bfseries,label={[font=\scriptsize,xshift=1mm]right:$BN$: $\frac{3}{10}$}] {N} edge from parent node[below,pos=.55,font=\scriptsize,text=teal!70!black,inner sep=1.5pt] {$\frac{2}{4}$}} edge from parent node[above,pos=.55,font=\scriptsize,text=teal!70!black,inner sep=1.5pt] {$\frac{3}{5}$}} child {node[circle,fill=black!80,text=white,inner sep=1.2pt,minimum size=4.2mm,font=\scriptsize\bfseries] {N} child {node[circle,fill=blue!70!black,text=white,inner sep=1.2pt,minimum size=4.2mm,font=\scriptsize\bfseries,label={[font=\scriptsize,xshift=1mm]right:$NB$: $\frac{3}{10}$}] {B} edge from parent node[above,pos=.55,font=\scriptsize,text=teal!70!black,inner sep=1.5pt] {$\frac{3}{4}$}} child {node[circle,fill=black!80,text=white,inner sep=1.2pt,minimum size=4.2mm,font=\scriptsize\bfseries,label={[font=\scriptsize,xshift=1mm]right:$NN$: $\frac{1}{10}$}] {N} edge from parent node[below,pos=.55,font=\scriptsize,text=teal!70!black,inner sep=1.5pt] {$\frac{1}{4}$}} edge from parent node[below,pos=.55,font=\scriptsize,text=teal!70!black,inner sep=1.5pt] {$\frac{2}{5}$}};
\end{tikzpicture}\end{center}
\begin{tasks}(1)
  \task $\frac35\cdot\frac24=\res{\frac3{10}}$.
  \task $\frac25\cdot\frac14=\res{\frac1{10}}$.
  \task $\frac3{10}+\frac3{10}=\res{\frac35}$. \ (Suman 1 \checkmark)
\end{tasks}

\ex{2} \par\smallskip\begin{center}\begin{tikzpicture}[grow'=right,level distance=19mm,level 1/.style={sibling distance=15mm},level 2/.style={sibling distance=7mm},every node/.style={font=\small\bfseries},edge from parent/.style={draw=muted,thick}]
\node[circle,fill=navy,inner sep=1.2pt] {} child {node[] {M} child {node[label={[font=\scriptsize,xshift=1mm]right:$0{,}45$}] {F} edge from parent node[above,pos=.55,font=\scriptsize,text=teal!70!black,inner sep=1.5pt] {$0{,}75$}} child {node[label={[font=\scriptsize,xshift=1mm]right:$0{,}15$}] {$\overline{F}$} edge from parent node[below,pos=.55,font=\scriptsize,text=teal!70!black,inner sep=1.5pt] {$0{,}25$}} edge from parent node[above,pos=.55,font=\scriptsize,text=teal!70!black,inner sep=1.5pt] {$0{,}6$}} child {node[] {$\overline{M}$} child {node[label={[font=\scriptsize,xshift=1mm]right:$0{,}1$}] {F} edge from parent node[above,pos=.55,font=\scriptsize,text=teal!70!black,inner sep=1.5pt] {$0{,}25$}} child {node[label={[font=\scriptsize,xshift=1mm]right:$0{,}3$}] {$\overline{F}$} edge from parent node[below,pos=.55,font=\scriptsize,text=teal!70!black,inner sep=1.5pt] {$0{,}75$}} edge from parent node[below,pos=.55,font=\scriptsize,text=teal!70!black,inner sep=1.5pt] {$0{,}4$}};
\end{tikzpicture}\end{center}$P(F/M)=0{,}75$ y $P(F/\overline{M})=0{,}25$ son probabilidades condicionadas.
\begin{tasks}(1)
  \task $0{,}6\cdot0{,}75=\res{0{,}45}$.
  \task $0{,}45+0{,}4\cdot0{,}25=0{,}45+0{,}1=\res{0{,}55}$.
  \task $0{,}4\cdot0{,}75=\res{0{,}3}$.
\end{tasks}

\ex{3} \par\smallskip\begin{center}\begin{tikzpicture}[grow'=right,level distance=18mm,level 1/.style={sibling distance=16mm},level 2/.style={sibling distance=8mm},every node/.style={font=\small\bfseries},edge from parent/.style={draw=muted,thick}]
\node[circle,fill=navy,inner sep=1.2pt] {} child {node[] {I} child {node[circle,fill=red!75!black,text=white,inner sep=1.2pt,minimum size=4.2mm,font=\scriptsize\bfseries,label={[font=\scriptsize,xshift=1mm]right:$\frac13\cdot\frac34=\frac14$}] {R} edge from parent node[above,pos=.55,font=\scriptsize,text=teal!70!black,inner sep=1.5pt] {$\frac34$}} child {node[circle,fill=blue!70!black,text=white,inner sep=1.2pt,minimum size=4.2mm,font=\scriptsize\bfseries] {B} edge from parent node[below,pos=.55,font=\scriptsize,text=teal!70!black,inner sep=1.5pt] {$\frac14$}} edge from parent node[above,pos=.55,font=\scriptsize,text=teal!70!black,inner sep=1.5pt] {$\frac13$}} child {node[] {II} child {node[circle,fill=red!75!black,text=white,inner sep=1.2pt,minimum size=4.2mm,font=\scriptsize\bfseries,label={[font=\scriptsize,xshift=1mm]right:$\frac23\cdot\frac14=\frac16$}] {R} edge from parent node[above,pos=.55,font=\scriptsize,text=teal!70!black,inner sep=1.5pt] {$\frac14$}} child {node[circle,fill=blue!70!black,text=white,inner sep=1.2pt,minimum size=4.2mm,font=\scriptsize\bfseries] {B} edge from parent node[below,pos=.55,font=\scriptsize,text=teal!70!black,inner sep=1.5pt] {$\frac34$}} edge from parent node[below,pos=.55,font=\scriptsize,text=teal!70!black,inner sep=1.5pt] {$\frac23$}};
\end{tikzpicture}\end{center}$P(R)=\frac13\cdot\frac34+\frac23\cdot\frac14=\frac14+\frac16=\res{\frac5{12}}$.

\ex{3} $\frac46\cdot\frac35\cdot\frac24=\frac{24}{120}=\res{\frac15}$.

\ex{3} Contrario: ninguna defectuosa, $P=\frac35\cdot\frac24=\frac3{10}$. \ $P(\text{al menos una})=1-\frac3{10}=\res{\frac7{10}}$.


% =====================================================
\seccio{7}{Tablas de contingencia}

\ex{1} 
\begin{tasks}(1)
  \task $\frac{8+6}{25}=\res{\frac{14}{25}}$.
  \task $\res{\frac6{25}}$.
  \task $\frac{8+4+6}{25}=\res{\frac{18}{25}}$.
  \task Solo las 12 chicas: $\frac8{12}=\res{\frac23}$.
\end{tasks}

\ex{2} \par\smallskip\begin{center}{\small\renewcommand{\arraystretch}{1.3}\begin{tabular}{l|cc|c}
\rowcolor{soft} & Federado & No federado & Total\\\hline
\cellcolor{soft}\textbf{Chicos} & 70 & 110 & \textbf{180}\\
\cellcolor{soft}\textbf{Chicas} & 50 & 170 & \textbf{220}\\\hline
\cellcolor{soft}\textbf{Total} & \textbf{120} & \textbf{280} & \textbf{400}\\
\end{tabular}}\end{center}
\begin{tasks}(1)
  \task $\frac{120}{400}=\res{\frac3{10}}$.
  \task $\frac{50}{400}=\res{\frac18}$.
  \task $\frac{50+170+70}{400}=\frac{290}{400}=\res{\frac{29}{40}}$.
  \task Solo contamos las 220 chicas: $\frac{50}{220}=\res{\frac5{22}}$.
\end{tasks}

\ex{2} \par\smallskip\begin{center}{\small\renewcommand{\arraystretch}{1.3}\begin{tabular}{l|cc|c}
\rowcolor{soft} & Gripe & Sin gripe & Total\\\hline
\cellcolor{soft}\textbf{Vacunados} & 12 & 228 & \textbf{240}\\
\cellcolor{soft}\textbf{No vacunados} & 48 & 112 & \textbf{160}\\\hline
\cellcolor{soft}\textbf{Total} & \textbf{60} & \textbf{340} & \textbf{400}\\
\end{tabular}}\end{center}Vacunados: $60\,\%$ de $400=240$, y con gripe, el $5\,\%$ de $240=12$. No vacunados: 160, y con gripe, el $30\,\%$ de $160=48$.
\begin{tasks}(1)
  \task $\frac{60}{400}=\res{0{,}15}$.
  \task $\frac{12}{240}=\res{0{,}05}$ y $\frac{48}{160}=\res{0{,}3}$.
  \task Solo las 60 personas con gripe: $\frac{12}{60}=\res{\frac15}$.
\end{tasks}

\ex{3} Hombres: $45\,\%$ de 200 $=90$; mujeres: 110; con gafas: $30\,\%$ de 200 $=60$.\par\smallskip\begin{center}{\small\renewcommand{\arraystretch}{1.3}\begin{tabular}{l|cc|c}
\rowcolor{soft} & Gafas & Sin gafas & Total\\\hline
\cellcolor{soft}\textbf{Hombres} & 20 & 70 & \textbf{90}\\
\cellcolor{soft}\textbf{Mujeres} & 40 & 70 & \textbf{110}\\\hline
\cellcolor{soft}\textbf{Total} & \textbf{60} & \textbf{140} & \textbf{200}\\
\end{tabular}}\end{center}
\begin{tasks}(1)
  \task $\frac{20}{200}=\res{\frac1{10}}$.
  \task $\frac{70}{200}=\res{\frac7{20}}$.
  \task Ahora los casos posibles son las 60 personas con gafas: $\frac{40}{60}=\res{\frac23}$.
\end{tasks}

\ex{3} \par\smallskip\begin{center}{\small\renewcommand{\arraystretch}{1.3}\begin{tabular}{l|cc|c}
\rowcolor{soft} & A pie & No a pie & Total\\\hline
\cellcolor{soft}\textbf{3.\textsuperscript{o}} & 12 & 18 & \textbf{30}\\
\cellcolor{soft}\textbf{4.\textsuperscript{o}} & 8 & 12 & \textbf{20}\\\hline
\cellcolor{soft}\textbf{Total} & \textbf{20} & \textbf{30} & \textbf{50}\\
\end{tabular}}\end{center}$P(\text{a pie})=\frac{20}{50}=\res{\frac25}$ \ y \ $P(\text{a pie}/\text{3.\textsuperscript{o}})=\frac{12}{30}=\res{\frac25}$. \ Saber que es de 3.\textsuperscript{o} no cambia la probabilidad: \rest{sí, son independientes}.

\ex{3} \par\smallskip\begin{center}{\small\renewcommand{\arraystretch}{1.3}\begin{tabular}{l|cc|c}
\rowcolor{soft} & Más de 50 € & 50 € o menos & Total\\\hline
\cellcolor{soft}\textbf{Tarjeta} & \cellcolor{accent!30}\textbf{140} & 560 & \textbf{700}\\
\cellcolor{soft}\textbf{Efectivo} & \cellcolor{accent!30}\textbf{30} & 270 & \textbf{300}\\\hline
\cellcolor{soft}\textbf{Total} & \textbf{170} & \textbf{830} & \textbf{1000}\\
\end{tabular}}\end{center}Tarjeta: 700, y 140 compran más de 50 €; efectivo: 300, y 30. En total, 170 clientes han comprado más de 50 €. \ $P(\text{tarjeta}/{>}50\text{ €})=\frac{140}{170}=\res{\frac{14}{17}}\approx0{,}82$.


\end{document}
